\section*{Capítulo \ref{ch:b1:atmospheric-circulation} --- La circulación general de la atmósfera}
% values: calc/ch28.py

\begin{solution}{exo:b1:atmospheric-circulation:1}
(a)~$f = 2\Omega\sin\varphi$: $0$ en el ecuador, \qty{7.29e-5}{s^{-1}} a \ang{30},
\qty{1.031e-4}{s^{-1}} a \ang{45} y \qty{1.458e-4}{s^{-1}} en el polo.
(b)~$fV = \num{1.031e-4} \times 10 = \qty{1.03e-3}{m.s^{-2}}$, unas \num{1.05e-4}
de $g$. (c)~En el hemisferio norte la desviación es a la derecha del movimiento: un viento que sopla hacia el norte es empujado hacia el este.
\end{solution}

\begin{solution}{exo:b1:atmospheric-circulation:2}
(a)~$|\nabla p|/\rho = (400/\num{4e5})/1.225 = \qty{8.16e-4}{m.s^{-2}}$.
(b)~$10/\num{8.16e-4} = \qty{12250}{s}$, \qty{3.4}{h}. (c)~La aceleración de Coriolis de un viento de \qty{10}{m/s} a \ang{45} N es \qty{1.03e-3}{m.s^{-2}}: las dos son del mismo tamaño (cociente $0.79$). En pocas horas desde que empieza a moverse, el aire siente una fuerza de Coriolis capaz de equilibrar el gradiente de presión: el viento se asienta cerca de la rapidez geostrófica, aquí $\num{8.16e-4}/\num{1.031e-4} = \qty{7.9}{m/s}$, en lugar de acelerar a lo largo del gradiente.
\end{solution}

\begin{solution}{exo:b1:atmospheric-circulation:3}
(a)~$f = 2 \times \num{7.292e-5} \times \sin\ang{40} = \qty{9.374e-5}{s^{-1}}$;
$|\vec v_g| = (400/\num{2.5e5})/(1.225 \times \num{9.374e-5}) = \qty{13.9}{m/s}$, o
$13.93/0.5144 = \qty{27.1}{kn}$. (b)~En el hemisferio norte el viento deja la baja a su izquierda: con la baja al sur, sopla hacia el oeste, desde el este (azimut \ang{90}). (c)~El asta apunta al este desde el círculo de la estación, hacia el viento; \qty{27.1}{kn} se redondea a \qty{25}{kn}, dos plumas enteras y media pluma, dibujadas en el lado horario del asta, aquí su lado sur, hacia la baja.
\end{solution}

\begin{solution}{exo:b1:atmospheric-circulation:4}
(a)~\ang{15} N: \omterm{def:b1:atmospheric-circulation:winds}{alisios} del noreste; \ang{15} S: \omterm{def:b1:atmospheric-circulation:winds}{alisios} del sureste; \ang{45} N y \ang{45} S: \omterm{def:b1:atmospheric-circulation:winds}{vientos del oeste}. (b)~A \ang{15} el aire superficial de la \omterm{def:b1:atmospheric-circulation:cells}{célula de Hadley} regresa hacia el ecuador y la fuerza de Coriolis lo gira hacia el oeste (a la derecha en el norte, a la izquierda en el sur): viento del este, $u$ negativo. A \ang{45} el aire superficial de la \omterm{def:b1:atmospheric-circulation:cells}{célula de Ferrel} se dirige hacia el polo y es girado hacia el este: viento del oeste, $u$ positivo. (c)~De las tres \omterm{def:b1:spherical-astronomy:geographic}{latitudes}, la presión es más alta a \ang{30} N, \qty{1016.8}{hPa}: la \omterm{def:b1:atmospheric-circulation:cells}{alta subtropical} yace cerca de \ang{30} N (el perfil completo sitúa su máximo cerca de \ang{35} N).
\end{solution}

\begin{solution}{exo:b1:atmospheric-circulation:5}
(a)~$|\nabla p|/\rho = (200/\num{5e5})/1.225 = \qty{3.27e-4}{m.s^{-2}}$, dividido por
$f$: \qty{2.6}{m/s} a \ang{60}, \qty{2.9}{m/s} a \ang{50}, \qty{4.5}{m/s} a
\ang{30}, \qty{6.5}{m/s} a \ang{20}, \qty{12.9}{m/s} a \ang{10} y
\qty{25.7}{m/s} a \ang{5}. (b)~La fuerza de Coriolis gira un viento en un tiempo del orden de
$1/f$: \qty{2.2}{h} a \ang{60}, \qty{11}{h} a \ang{10}, \qty{22}{h} a \ang{5}, e infinito en el ecuador. Cerca del ecuador el campo de presión cambia, y la fricción actúa, mucho antes de que el viento pueda girarse a lo largo de las isobaras: el equilibrio nunca se establece, y el aire fluye en parte a lo largo del gradiente, a velocidades moderadas.
\end{solution}

\begin{solution}{exo:b1:atmospheric-circulation:6}
(a)~$u = 464.6\,\sin^2\varphi/\cos\varphi$: \qty{14.2}{m/s} a \ang{10},
\qty{32.2}{m/s} a \ang{15}, \qty{57.8}{m/s} a \ang{20}, \qty{134.1}{m/s} a
\ang{30}. (b)~$u = 464.6\,(\cos^2\ang{10} - \cos^2\varphi)/\cos\varphi$:
\qty{17.7}{m/s} a \ang{15}, \qty{42.9}{m/s} a \ang{20}, \qty{117.9}{m/s} a
\ang{30}. (c)~El aire que parte más lejos del ecuador parte más cerca del eje, con menos momento angular de sobra: llega con un viento del oeste más débil a cada \omterm{def:b1:spherical-astronomy:geographic}{latitud}. Extender la rama ascendente baja el chorro predicho, pero no lo baja lo suficiente para alcanzar los \qty{40}{m/s} observados.
\end{solution}

\begin{solution}{exo:b1:atmospheric-circulation:7}
(a)~El viento sopla hacia el noreste. Con la espalda al viento, mirando al noreste, la baja queda a la izquierda: al noroeste. (b)~El viento sopla hacia el sureste; en el hemisferio sur la baja queda a la derecha de quien tiene la espalda al viento: al suroeste. (c)~El viento superficial se desvía de las isobaras hacia la baja, es decir, en sentido antihorario respecto al \omterm{def:b1:atmospheric-circulation:geostrophic}{viento geostrófico} en el hemisferio norte. El \omterm{def:b1:atmospheric-circulation:geostrophic}{viento geostrófico}, y la baja a ángulos rectos a su izquierda, quedan por tanto girados en sentido horario respecto a la respuesta de (a): la baja yace entre el noroeste y el norte, en el azimut $\ang{315} + \alpha$ para un ángulo de cruce
$\alpha$.
\end{solution}

\begin{solution}{exo:b1:atmospheric-circulation:8}
(a)~El área entre el ecuador y \ang{30} N es $2\pi a_\oplus^2\sin\ang{30} =
\qty{1.275e14}{m^2}$; la \omterm{def:b1:observing-the-sky:extinction}{masa de aire} por metro cuadrado entre los dos niveles es
$\Delta p/g = \num{9e4}/9.81 = \qty{9174}{kg}$; el total es
\qty{1.170e18}{kg}. (b)~$\num{1.170e18}/\num{8.7e10} = \qty{1.34e7}{s}$, unos
\num{156} días, cinco meses. (c)~$\num{1.170e18}/\num{21.4e10} = \qty{5.47e6}{s}$,
unos \num{63} días, dos meses.
\end{solution}

\begin{solution}{exo:b1:atmospheric-circulation:9}
(a)~$\num{0.01} \times \num{86400} = \qty{864}{m}$ al día. (b)~$0.864 \times 9.8 =
\qty{8.5}{K}$ al día. (c)~El aire se enfría emitiendo radiación infrarroja: la atmósfera pierde más energía por radiación de la que absorbe, y el descenso del aire, que lo calienta por compresión, compensa la pérdida. (d)~El aire que desciende se calienta y se seca respecto a la saturación, sus nubes se evaporan y la lluvia no puede formarse: las \omterm{def:b1:atmospheric-circulation:cells}{altas subtropicales} son cinturones de cielos despejados y poca lluvia, donde se forman desiertos cálidos en tierra.
\end{solution}

\begin{solution}{exo:b1:atmospheric-circulation:10}
(a)~$V^2 + frV - fr|\vec v_g| = 0$ con $fr = \qty{55.85}{m/s}$:
$V = \tfrac12\bigl(-55.85 + \sqrt{55.85^2 + 4 \times 55.85 \times 15}\bigr) =
\qty{12.3}{m/s}$, \qty{82}{\percent} del \omterm{def:b1:atmospheric-circulation:geostrophic}{viento geostrófico}. (b)~$V^2 - frV +
fr|\vec v_g| = 0$ tiene raíces reales solo si $(fr)^2 \geq 4fr|\vec v_g|$, es decir,
$|\vec v_g| \leq fr/4 = \qty{14.0}{m/s}$. (c)~Cerca del centro de un alto, donde $r$ es
pequeño, el límite $fr/4$ es pequeño: el gradiente de presión debe ser débil, las isobaras muy separadas y los vientos suaves. Una depresión no tiene tal límite: sus isobaras pueden agolparse cerca del centro, con vientos fuertes.
\end{solution}

\begin{solution}{exo:b1:atmospheric-circulation:11}
(a)~Para $\varphi$ pequeño, $\sin\varphi \approx \varphi$ y $\cos\varphi \approx 1$,
así que $u \approx \Omega a_\oplus\varphi^2$. (b)~$\ang{10} = \qty{0.1745}{rad}$:
$464.6 \times 0.1745^2 = \qty{14.15}{m/s}$, frente a \qty{14.23}{m/s} exacto,
\qty{0.5}{\percent} menos. (c)~Resolviendo $464.6\sin^2\varphi/\cos\varphi = 27.0$
numéricamente se obtiene $\varphi = \ang{13.7}$. (d)~El anillo ya coincide con la \omterm{def:b1:atmospheric-circulation:jet}{corriente en chorro} real anual media a menos de \ang{14}: la conservación del momento angular describe la rama superior de la \omterm{def:b1:atmospheric-circulation:cells}{célula de Hadley} solo en los trópicos profundos; hacia el polo, desde unos \ang{15}, los remolinos toman el mando.
\end{solution}

\begin{solution}{exo:b1:atmospheric-circulation:12}
(a)~Media $31.7/12 = \ang{2.6}$ N; intervalo desde \ang{12.8} S hasta \ang{10.9} N, es decir,
\ang{23.7}. (b)~La posición más septentrional llega en septiembre, tres meses después del solsticio de junio; la más meridional en febrero, dos meses después del solsticio de diciembre. El retraso procede del almacenamiento de calor, sobre todo en la \omterm{def:b1:fluid-envelopes-structure:ocean-layers}{capa de mezcla} del océano (\cref{ch:b1:fluid-envelopes-structure}), que se calienta y se enfría despacio. (c)~La media zonal promedia todas las longitudes: si sobre los continentes la banda de lluvia se desplaza mucho hacia el hemisferio de verano, la media de todas las longitudes puede cruzar el ecuador aunque la banda sobre un océano no lo haga.
\end{solution}

\begin{solution}{pb:b1:atmospheric-circulation:1}
\begin{enumerate}
  \item $\rho = p/(R_dT) = \num{100600}/(287.05 \times 283.15) =
        \qty{1.238}{kg.m^{-3}}$.
  \item $|\nabla p| = \qty{400}{Pa}/\qty{2.5e5}{m} = \qty{1.6e-3}{Pa.m^{-1}}$.
  \item $|\nabla p|/\rho = \num{1.6e-3}/1.238 = \qty{1.293e-3}{m.s^{-2}}$.
  \item $\Omega = 2\pi/\num{86164.2} = \qty{7.292e-5}{rad.s^{-1}}$;
        $f = 2\Omega\sin\ang{50} = \qty{1.117e-4}{s^{-1}}$.
  \item $|\vec v_g| = \num{1.293e-3}/\num{1.117e-4} = \qty{11.57}{m/s}$.
  \item $11.57/0.514444 = \qty{22.5}{kn}$; $11.57 \times 3.6 = \qty{41.7}{km/h}$.
  \item El tiempo es $|\vec v_g|/(|\nabla p|/\rho)$ y $|\vec v_g| = (|\nabla
        p|/\rho)/f$, así que es $1/f = \qty{8951}{s}$, \qty{2.49}{h}; y $2\pi/f =
        \qty{15.6}{h}$. Los patrones que duran un día o más están, por tanto, cerca del equilibrio.
  \item El gradiente de presión apunta hacia el norte, hacia la baja. El \omterm{def:b1:atmospheric-circulation:geostrophic}{viento geostrófico} sopla a lo largo de las isobaras con la baja a su izquierda: hacia el este, desde el oeste, azimut \ang{270}: un viento del oeste.
  \item Con la espalda al viento del oeste, mirando al este, la izquierda es el norte, donde está la baja: la ley se cumple.
  \item El asta apunta al oeste desde el círculo de la estación, hacia de donde viene el viento; \qty{22.5}{kn} se redondea a \qty{20}{kn}, dos plumas enteras; están en el lado horario del asta, su lado norte, hacia la baja.
  \item A \ang{50} S, $f < 0$: el viento deja la baja a su derecha, así que sopla hacia el oeste, desde el este.
  \item $f(\ang{30}) = \qty{7.292e-5}{s^{-1}}$: \qty{17.7}{m/s}; $f(\ang{60}) =
        \qty{1.263e-4}{s^{-1}}$: \qty{10.2}{m/s}.
  \item $f(\ang{5}) = \qty{1.271e-5}{s^{-1}}$ da \qty{101.7}{m/s}, un viento nunca observado con gradientes así. El tiempo de giro $1/f$ es \qty{22}{h}:
        el equilibrio no tiene tiempo de establecerse, y la fórmula no se aplica.
  \item $9.0/11.57 = 0.778$; a lo largo de las isobaras $9.0\cos\ang{25} = \qty{8.16}{m/s}$,
        a través de ellas hacia la baja $9.0\sin\ang{25} = \qty{3.80}{m/s}$.
  \item El \omterm{def:b1:atmospheric-circulation:geostrophic}{viento geostrófico} sopla hacia el azimut \ang{90}; girado \ang{25} hacia la baja en el norte, el viento de superficie sopla hacia \ang{65}, desde
        \ang{245}, el oeste-suroeste.
  \item $\tan\alpha = 1/(f\tau)$ da $\tau = 1/(\num{1.117e-4}\tan\ang{25}) =
        \qty{1.92e4}{s}$, \qty{5.3}{h}; el modelo predice $11.57\cos\ang{25} =
        \qty{10.5}{m/s}$, más que los \qty{9.0}{m/s} observados: la fricción real no es exactamente lineal en el viento, y frena más que el modelo en este ángulo.
  \item A través de un cilindro de radio $R = \qty{500}{km}$ y altura $h =
        \qty{1}{km}$: $\dot m = 2\pi R h\rho u_r = 2\pi \times \num{5e5} \times 1000
        \times 1.238 \times 3.80 = \qty{1.48e10}{kg/s}$.
  \item $w = \dot m/(\rho\pi R^2) = 2hu_r/R = 2 \times 1000 \times 3.80/\num{5e5} =
        \qty{0.0152}{m/s}$, unos \qty{1310}{m} al día.
  \item $1.31 \times 9.8 = \qty{12.9}{K}$ al día: el aire pronto alcanza la saturación, y la depresión se llena de nubes y lluvia. Bajo un \omterm{def:b1:atmospheric-circulation:cyclone}{anticiclón} el aire de superficie fluye hacia fuera, el aire de arriba desciende y se calienta, y el cielo se despeja.
  \item $\Omega a_\oplus = \num{7.292e-5} \times \num{6.371e6} = \qty{464.6}{m/s}$
        en el ecuador; $464.6\cos\ang{30} = \qty{402.3}{m/s}$ a \ang{30} N.
  \item $u = 464.6\,\sin^2\varphi/\cos\varphi$: \qty{57.8}{m/s} a \ang{20} N y
        \qty{134.1}{m/s} a \ang{30} N.
  \item $u = 464.6\,(\cos^2\ang{10} - \cos^2\ang{30})/\cos\ang{30} =
        \qty{117.9}{m/s}$.
  \item $134.1/42.3 = 3.17$. Los remolinos de las \omterm{def:b1:spherical-astronomy:geographic}{latitudes} medias, depresiones y \omterm{def:b1:atmospheric-circulation:cyclone}{anticiclones}, toman momento angular de la rama superior y lo llevan a otra parte; y el aire de la célula sube sobre una banda de \omterm{def:b1:spherical-astronomy:geographic}{latitudes} en lugar de estar en reposo sobre el ecuador. La media zonal además promedia un chorro serpenteante sobre todas las longitudes.
  \item $M/M_0 = (402.3 + 42.3)\cos\ang{30}/464.6 = 0.829$: el aire ha perdido
        \qty{17}{\percent} del momento angular que habría conservado como anillo libre.
  \item El \omterm{def:b1:atmospheric-circulation:geostrophic}{viento geostrófico} en P es \qty{11.6}{m/s}, desde el oeste; el viento medido en el suelo es \qty{9.0}{m/s}, unas cuatro quintas partes de él, y girado \ang{25} hacia la baja por la fricción.
\end{enumerate}
\end{solution}
